\subsection{Purpose}

YASF is a deterministic specification language used in the current CGD implementation as the form of interface-level intent representation $I$ in the architectural pair $D = \langle I, R \rangle$. YASF describes a system in controlled natural language with a canonical set of sections. A YASF text is compiled into the structural realization $R = \langle G, \Sigma, C \rangle$, where $G$ is a typed graph of functions and tables, $\Sigma$ is a set of table signatures, and $C$ is a set of function contracts. YASF plays the role of the authoring layer: it is the format in which a person fixes intent before it is translated into the formal core.

\subsection{Design principles}

\textbf{Principle 1 --- controlled natural language.} YASF is written in natural language, but with a fixed section structure and controlled terminology. This yields a low entry threshold while preserving compilability.

\textbf{Principle 2 --- deterministic compilation.} For the same consistent YASF text, the compiler produces the same structural realization $R = \langle G, \Sigma, C \rangle$ with the same graph relations, contracts, and signatures. Determinism is treated as a design invariant and is verified through conformance tests.

\textbf{Principle 3 --- minimality.} YASF describes exactly what is required to build the structural realization: functions, tables, data flows, constraints, and contracts. Implementation details such as algorithms, data structures, and coding style are not part of YASF.

\textbf{Principle 4 --- modularity.} Each function is described by a separate YASF block. Blocks may refer to one another through function names and table names. This enables independent development, verification, and compilation of separate parts of intent.

\textbf{Principle 5 --- traceability.} Every element of the structural realization---a node or edge in $G$, a signature $\Sigma(t)$, or a contract $C(f)$---must be traceable to a specific YASF section. Conversely, each YASF section must produce specific elements in $R = \langle G, \Sigma, C \rangle$.

\subsection{Structure of a YASF block}

A YASF block describes one function. Its canonical structure includes the following sections.

\textbf{Purpose} --- a short natural-language description of the function (1--2 sentences). Compiles into: a \texttt{function} node in graph $G$ with a name and description.

\textbf{Input data} --- the list of tables consumed by the function, indicating fields, types, keys, and \texttt{kind}. Compiles into: \texttt{table} nodes in $G$, \texttt{consumes} edges, the $C_{\mathrm{in}}$ component of the contract, and signatures $\Sigma(t)$ of input tables.

\textbf{Reference data} --- the list of tables read by the function for reference. Compiles into: \texttt{reads} edges, the $R$ component of the contract, and---when needed---the signatures of the corresponding tables.

\textbf{Processing} --- the sequence of processing steps in controlled natural language. The order of items gives the basic sequence of steps. The concrete graph edges and contract dependencies are formed during compilation jointly with table signatures and the other sections of the block.

\textbf{Output data} --- the list of tables produced by the function, indicating fields, types, keys, and \texttt{kind}. Compiles into: \texttt{produces} edges, the $C_{\mathrm{out}}$ component of the contract, and signatures $\Sigma(t)$ of output tables.

\textbf{Errors} --- the list of error tables produced by the function when it deviates from the success path. Compiles into: \texttt{produces} edges, additional $C_{\mathrm{out}}$ elements, and signatures $\Sigma(t)$ of error tables (\texttt{kind: error}).

\textbf{Side effects} --- the list of tables written by the function that are not its main outputs. Compiles into: \texttt{writes} edges, the $W$ component of the contract, and when needed the signatures of the corresponding tables.

\textbf{Triggers} --- the list of functions activated by the given function on completion. Compiles into: \texttt{triggers} edges and the $T_{\mathrm{out}}$ component of the contract.

\textbf{Contract} --- an explicit fixation of the six-component contract $C(f) = \langle C_{\mathrm{in}}, C_{\mathrm{out}}, R, W, T_{\mathrm{in}}, T_{\mathrm{out}} \rangle$. It may be written manually by the author or generated by the compiler from the other sections. If both are present, the compiler checks consistency.

\subsection{Compilation rules: YASF $\to$ $R = \langle G, \Sigma, C \rangle$}

\textbf{Rule 1 --- node extraction.} Every function name mentioned in the ``Purpose'' or ``Triggers'' sections produces a node with $\tau_V = \texttt{function}$. Every table name mentioned in ``Input data,'' ``Reference data,'' ``Output data,'' ``Errors,'' or ``Side effects'' produces a node with $\tau_V = \texttt{table}$.

\textbf{Rule 2 --- edge extraction.} A table from ``Input data'' produces a \texttt{consumes} (table $\to$ function) edge. A table from ``Output data'' or ``Errors'' produces a \texttt{produces} (function $\to$ table) edge. A table from ``Reference data'' produces a \texttt{reads} (table $\to$ function) edge. A table from ``Side effects'' produces a \texttt{writes} (function $\to$ table) edge. A function from ``Triggers'' produces a \texttt{triggers} (function $\to$ function) edge.

\textbf{Rule 3 --- signature extraction.} For every table described with fields, types, and keys, a signature $\Sigma(t) = \langle \mathrm{Fields}, \mathrm{Keys}, \mathrm{Kind}, \mathrm{Constraints} \rangle$ is formed according to the definition of the paper's formal core.

\textbf{Rule 4 --- contract formation.} The contract $C(f)$ is formed from the extracted edges. If the author explicitly specified a contract in the ``Contract'' section, the compiler checks equality between the generated and author-specified contract. A mismatch is treated as a structural defect.

\textbf{Rule 5 --- completeness checking.} After a block is compiled, the following is checked: all tables appearing in the contract have defined signatures, all edges are type-correct, no nodes are left without edges, and the resulting local structural realization is consistent with the verification profile of the current step.

\subsection{Mapping between YASF sections and contract components}

\begin{tabular}{llll}
\hline
\textbf{YASF section} & \textbf{Contract} & \textbf{Edge type} & \textbf{Direction} \\
\hline
Input data & $C_{\mathrm{in}}$ & consumes & table $\to$ function \\
Reference data & $R$ & reads & table $\to$ function \\
Output data & $C_{\mathrm{out}}$ & produces & function $\to$ table \\
Errors & $C_{\mathrm{out}}$ & produces & function $\to$ table \\
Side effects & $W$ & writes & function $\to$ table \\
Triggers & $T_{\mathrm{out}}$ & triggers & function $\to$ function \\
External block & $T_{\mathrm{in}}$ & triggers & function $\to$ function \\
\hline
\end{tabular}

\vspace{0.5em}
The $T_{\mathrm{in}}$ component has no dedicated section inside a function block: it is derived from the ``Triggers'' sections of other blocks that refer to this function.

\subsection{Example of a YASF block}

\begin{lstlisting}
Function: RegisterUser
Purpose: Register a new user by email.

Input data:
RegistrationRequest (id PK, name: string required,
  email: string required, timestamp: timestamp required;
  kind: event)

Reference data:
User (id PK, name: string, email: string, status: string,
  created_at: timestamp; kind: entity)
  -- read to check email uniqueness.

Processing:
1. Accept RegistrationRequest.
2. Check format and completeness of the data.
3. Check email uniqueness in the User table.
4. Create a User record with status pending.
5. Record RegistrationEvent.
6. Write the action to AuditLog.
7. Trigger SendConfirmation.

Output data:
User (id PK, name: string, email: string, status: string,
  created_at: timestamp; kind: entity)
RegistrationEvent (id PK, user_id FK->User,
  event_type: string, timestamp: timestamp; kind: event)

Errors:
RegistrationError (id PK, error_code: string,
  message: string, timestamp: timestamp; kind: error)
  -- for invalid input or duplicate email.

Side effects:
AuditLog (id PK, action: string, entity_id: string,
  timestamp: timestamp; kind: log)

Triggers:
SendConfirmation
  -- send registration confirmation.

Contract v2:
C(RegisterUser) = <Cin: {RegistrationRequest},
  Cout: {User, RegistrationEvent, RegistrationError},
  R: {User}, W: {AuditLog},
  Tin: empty, Tout: {SendConfirmation}>
\end{lstlisting}

\subsection{Decomposition in YASF}

When a function is decomposed into subfunctions, each subfunction receives its own YASF block. The original block is preserved as an aggregated representation with a note such as:

\begin{lstlisting}
Function: RegisterUser (aggregated node)
Decomposition: ValidateInput -> CheckUniqueness ->
  CreateUser -> (triggers) SendConfirmation
Refinement map: see blocks ValidateInput,
  CheckUniqueness, CreateUser, SendConfirmation
\end{lstlisting}

Each subfunction block has the full structure from \S F.3. Intermediate tables (\texttt{ValidatedRequest}, \texttt{UniqueRequest}) are described in the ``Output data'' and ``Input data'' sections of the corresponding subfunctions. Formal computation of the external contract of the subgraph in such a decomposition is moved to Appendix~\ref{app:subgraph-contract}.

\subsection{Verification-driven dialogue in YASF}

When the verifier discovers defects after compiling a YASF block, they are fixed and classified as follows.

\textbf{Blocking} --- require the author to extend the YASF text. Example: the ``Errors'' section is absent or empty while the ``Processing'' section contains branching conditions. The verifier generates a question; after the author answers, the section is extended.

\textbf{Auto-resolvable} --- the compiler or verifier can extend the block automatically. Example: if an error path has been added and the function already has an \texttt{AuditLog} side effect, an \texttt{AuditLog} entry on the error path is added by analogy.

\textbf{Advisory} --- warnings recorded as comments on the YASF block. Example: a table from ``Output data'' is not consumed by any known function.

After all blocking defects have been resolved and auto-resolvable fixes have been applied, the block is recompiled into $R = \langle G, \Sigma, C \rangle$ and checked again until specification closure is reached (see \S 5 of the main text).

\subsection{Conformance tests}

Compilation determinism is verified through conformance tests---a fixed set of pairs:

\begin{tabular}{ll}
\hline
\textbf{Input} & \textbf{Expected output} \\
\hline
YASF text (function block) & Subgraph $G$: nodes, edges, types \\
YASF text (function block) & Contract $C(f)$: six components \\
YASF text (table description) & Signature $\Sigma(t)$: Fields, Keys, Kind, Constraints \\
YASF text (full block) & Local structural realization $R = \langle G, \Sigma, C \rangle$ \\
\hline
\end{tabular}

\vspace{0.5em}
A test is considered passed if the compiler produces a structurally identical local realization, up to ordering of elements in sets. The conformance-test set is extended as new YASF constructs appear.

\subsection{Limitations of the current YASF version}

Not supported: a formal grammar (YASF is defined through section structure and compilation rules rather than through BNF/EBNF), automatic semantic processing of free text inside the ``Processing'' section (the compiler relies on mentioned function and table names rather than on the full semantics of free text), inheritance between blocks (one YASF block cannot extend another), and conditional sections (branching inside a block is not formalized; the error path is described through a separate ``Errors'' section).

Planned: a formal grammar for YASF (to enable automatic syntax validation), conditional sections (to describe explicit branching inside a block), and template support (typical block patterns that can be inherited and overridden).

\subsection{Cross-references}

\begin{tabular}{ll}
\hline
\textbf{From Appendix F} & \textbf{Reference} \\
\hline
Rules 1--2 (node and edge extraction) & \S 4.1 (graph $G$ and edge types) \\
Rule 3 (signature extraction) & \S 4.2 (table signatures $\Sigma$) \\
Rule 4 (contract formation) & \S 4.3 (function contracts $C$) \\
Rule 5 (completeness checking) & \S 5.4 (specification closure) \\
Decomposition & Appendix C (ComputeSubgraphContract) \\
Verification-driven dialogue & \S\S 5.3--5.4 \\
Conformance tests & \S 3.1 (YASF as the intent interface) \\
\hline
\end{tabular}
